\documentclass[14pt, a4paper]{extarticle} 

% Поддержка русского языка и кодировок
\usepackage[utf8]{inputenc}
\usepackage[T2A]{fontenc}
\usepackage[russian]{babel}

% Настройка полей (ГОСТ)
\usepackage[left=3cm, right=1.5cm, top=2cm, bottom=2cm]{geometry}

% Графика и цвета
\usepackage{graphicx}
\usepackage{xcolor}

% Таблицы и списки
\usepackage{tabularx}
\usepackage{booktabs}
\usepackage{longtable}
\usepackage{enumitem}
\setlist{nosep}

% Математика
\usepackage{amsmath}
\usepackage{amssymb}
\usepackage{amsfonts}

% Ссылки и метаданные
\usepackage{hyperref}
\hypersetup{
    colorlinks=true,
    linkcolor=blue,
    filecolor=magenta,      
    urlcolor=blue,
    citecolor=black,
    pdftitle={Описание программных модулей системы APS},
}

% Вставка исходного кода
\usepackage{listings}
\definecolor{codegreen}{rgb}{0,0.6,0}
\definecolor{codegray}{rgb}{0.5,0.5,0.5}
\definecolor{codepurple}{rgb}{0.58,0,0.82}
\definecolor{backcolour}{rgb}{0.96,0.96,0.96}

\lstdefinestyle{mystyle}{
    backgroundcolor=\color{backcolour},   
    commentstyle=\color{codegreen}\itshape,
    keywordstyle=\color{blue}\bfseries,
    numberstyle=\tiny\color{codegray},
    stringstyle=\color{codepurple},
    basicstyle=\ttfamily\footnotesize,
    breakatwhitespace=false,         
    breaklines=true,                 
    captionpos=b,                    
    keepspaces=true,                 
    numbers=left,                    
    numbersep=5pt,                  
    showspaces=false,                
    showstringspaces=false,
    showtabs=false,                  
    tabsize=2,
    frame=single,
    rulecolor=\color{black!30},
    extendedchars=true,
    inputencoding=utf8,
    literate={а}{{\cyra}}1 {б}{{\cyrb}}1 {в}{{\cyrv}}1 {г}{{\cyrg}}1 {д}{{\cyrd}}1 
             {е}{{\cyre}}1 {ё}{{\cyryo}}1 {ж}{{\cyrzh}}1 {з}{{\cyrz}}1 {и}{{\cyri}}1 
             {й}{{\cyrishrt}}1 {к}{{\cyrk}}1 {л}{{\cyrl}}1 {м}{{\cyrm}}1 {н}{{\cyrn}}1 
             {о}{{\cyro}}1 {п}{{\cyrp}}1 {р}{{\cyrr}}1 {с}{{\cyrs}}1 {т}{{\cyrt}}1 
             {у}{{\cyru}}1 {ф}{{\cyrf}}1 {х}{{\cyrh}}1 {ц}{{\cyrc}}1 {ч}{{\cyrch}}1 
             {ш}{{\cyrsh}}1 {щ}{{\cyrshch}}1 {ъ}{{\cyrhrdsn}}1 {ы}{{\cyrery}}1 
             {ь}{{\cyrsftsn}}1 {э}{{\cyrerev}}1 {ю}{{\cyryu}}1 {я}{{\cyrya}}1
             {А}{{\CYRA}}1 {Б}{{\CYRB}}1 {В}{{\CYRV}}1 {Г}{{\CYRG}}1 {Д}{{\CYRD}}1 
             {Е}{{\CYRE}}1 {Ё}{{\CYRYO}}1 {Ж}{{\CYRZH}}1 {З}{{\CYRZ}}1 {И}{{\CYRI}}1 
             {Й}{{\CYRISHRT}}1 {К}{{\CYRK}}1 {Л}{{\CYRL}}1 {М}{{\CYRM}}1 {Н}{{\CYRN}}1 
             {О}{{\CYRO}}1 {П}{{\CYRP}}1 {Р}{{\CYRR}}1 {С}{{\CYRS}}1 {Т}{{\CYRT}}1 
             {У}{{\CYRU}}1 {Ф}{{\CYRF}}1 {Х}{{\CYRH}}1 {Ц}{{\CYRC}}1 {Ч}{{\CYRCH}}1 
             {Ш}{{\CYRSH}}1 {Щ}{{\CYRSHCH}}1 {Ъ}{{\CYRHRDSN}}1 {Ы}{{\CYRERY}}1 
             {Ь}{{\CYRSFTSN}}1 {Э}{{\CYREREV}}1 {Ю}{{\CYRYU}}1 {Я}{{\CYRYA}}1
             {°}{{$^\circ$}}1
}

\lstset{style=mystyle}

\begin{document}

% ========================================================================
% ТИТУЛЬНЫЙ ЛИСТ
% ========================================================================
\begin{titlepage}
\begin{center}
\vspace*{2cm}

\textbf{\Large ОПИСАНИЕ ПРОГРАММНЫХ МОДУЛЕЙ}\\[0.5cm]
\textbf{\Large СИСТЕМЫ ПЛАНИРОВАНИЯ ПРОИЗВОДСТВА (APS)}\\[0.5cm]
\textbf{\Large НА ОСНОВЕ ЦИФРОВОГО ДВОЙНИКА}\\[0.5cm]
\textbf{\Large ГАЛЬВАНИЧЕСКОГО КОМПЛЕКСА}\\[2cm]

{\large Приложение к диссертационной работе}\\[3cm]

\vfill

{\large 2026}
\end{center}
\end{titlepage}

% ========================================================================
\tableofcontents
\newpage

% ========================================================================
\section{Введение}
% ========================================================================

Настоящий документ содержит подробное описание каждого программного модуля (скрипта и класса), входящего в состав распределённой системы автоматизированного динамического планирования производства (APS --- Advanced Planning and Scheduling) для технологического комплекса горячего цинкования. 

Система представляет собой цифровой двойник промышленной гальванической линии, состоящей из 7 последовательных технологических стадий, и включает следующие взаимосвязанные подсистемы:

\begin{enumerate}
    \item \textbf{Серверное ядро} (Spring Boot, Java 17+) --- центральный узел вычислений алгоритма APS, встроенный сервер промышленного протокола OPC UA (на базе Eclipse Milo), HTTP REST API для веб- и мобильных клиентов;
    \item \textbf{Аппаратный мост и телеметрия} (FastAPI, Python 3.10+, PySerial) --- аппаратно-программный шлюз Hardware-in-the-Loop (HIL), выполняющий чтение телеметрии физического микроконтроллера ESP32 (датчики DHT11/DHT22) по последовательному интерфейсу UART (COM-порт 115200 бод), локальную симуляцию смежных стадий и ретрансляцию данных в Java-сервер;
    \item \textbf{Визуальный 2D-симулятор} (PyGame, Python, AsyncUA) --- графический SCADA-клиент OPC UA, визуализирующий в реальном времени работу многопозиционной технологической линии, физику перемещения манипулятора (мостового крана) и динамику технологических параметров ванн;
    \item \textbf{Единый оркестратор запуска} (Launcher, C\# / .NET) --- системная утилита координации жизненного цикла процессов, проверки доступности портов, автоматического запуска веб-браузера и безопасного аварийного завершения дочерних процессов;
    \item \textbf{Веб-дашборд} (React 18 / Vite, TypeScript, Tailwind CSS) --- интерактивная панель мониторинга цифрового двойника, отображающая ход выполнения блок-схемы алгоритма APS, телеметрию всех 7 технологических стадий, очереди MES и системные логи;
    \item \textbf{Мобильный терминал} (Android, Kotlin / Jetpack Compose) --- мобильный SCADA-терминал оператора линии для оперативного удалённого мониторинга и ручной корректировки очереди.
\end{enumerate}

Все программные модули описаны в строгой технической форме с приведением архитектурных особенностей, алгоритмических формул и ключевых фрагментов исходного кода (листингов).

% ========================================================================
\section{Архитектура системы и протоколы взаимодействия}
% ========================================================================

Система построена на принципах сервис-ориентированной распределённой архитектуры. Обмен данными между компонентами осуществляется по четырём специализированным каналам связи:

\begin{itemize}
    \item \textbf{UART / Serial} (COM20, 115200 бод, 8N1) --- физический канал связи с контроллером ESP32 для сбора первичных показаний температуры и влажности сушильной камеры;
    \item \textbf{OPC UA} (TCP-порт 4840) --- бинарный промышленный протокол связи реального времени (спецификация IEC 62541). Обеспечивает циклическую передачу телеметрии и команд управления между Java-сервером и PyGame-симулятором с частотой 5--10~Гц;
    \item \textbf{REST API Java-сервера} (HTTP-порт 8080) --- интерфейс интеграции веб-дашборда, Android-приложения и аппаратного моста. Поддерживает поллинг с частотой 2~Гц;
    \item \textbf{FastAPI Аппаратного моста} (HTTP-порт 8000) --- локальный REST-сервер для автономного мониторинга телеметрии датчиков и ретрансляции данных.
\end{itemize}

\begin{table}[h!]
\centering
\caption{Пакетная структура серверного ядра (Java)}
\begin{tabularx}{\textwidth}{|l|X|}
\hline
\textbf{Пакет} & \textbf{Назначение} \\
\hline
\texttt{com.smartfactory} & Точка входа в приложение и конфигурация планировщика задач \\
\texttt{com.smartfactory.model} & Модели данных (партии \texttt{BatchInfo}, параметры комплекса \texttt{ComplexParameters}, суперпозиция \texttt{SuperpositionData}, действия \texttt{ActionSequence}) \\
\texttt{com.smartfactory.state} & Глобальное состояние системы и 7 технологических стадий (\texttt{ApiState}) \\
\texttt{com.smartfactory.service} & Бизнес-логика алгоритма APS, управление расписанием, симуляция процессов \\
\texttt{com.smartfactory.controller} & Контроллер REST API для клиентов и аппаратного моста \\
\texttt{com.smartfactory.opcua} & Встроенный OPC UA сервер и адресное пространство узлов \\
\hline
\end{tabularx}
\end{table}

% ========================================================================
\newpage
\section{Серверное ядро (Java / Spring Boot)}
% ========================================================================

% ------------------------------------------------------------------------
\subsection{SmartFactoryApplication.java --- Точка входа приложения}
% ------------------------------------------------------------------------

\textbf{Путь:} \texttt{com.smartfactory.SmartFactoryApplication}\\
\textbf{Зависимости:} Spring Boot 3.x, Spring Scheduling, Eclipse Milo

Класс инициализирует контекст Spring Boot и настраивает пул потоков планировщика задач. По умолчанию стандартный планировщик Spring Boot использует один поток, что создаёт риск взаимных блокировок между высокочастотным опросом OPC UA, вычислениями алгоритма APS и симуляцией датчиков. 

Для исключения взаимного влияния настраивается специализированный бин \texttt{TaskScheduler} с фиксированным пулом из 5 независимых потоков (\texttt{aps-scheduler-}).

\begin{lstlisting}[language=Java, caption=Фрагмент SmartFactoryApplication.java]
package com.smartfactory;

import org.springframework.boot.SpringApplication;
import org.springframework.boot.autoconfigure.SpringBootApplication;
import org.springframework.context.annotation.Bean;
import org.springframework.scheduling.TaskScheduler;
import org.springframework.scheduling.annotation.EnableScheduling;
import org.springframework.scheduling.concurrent.ThreadPoolTaskScheduler;

@SpringBootApplication
@EnableScheduling
public class SmartFactoryApplication {

    public static void main(String[] args) {
        SpringApplication.run(SmartFactoryApplication.class, args);
    }

    @Bean
    public TaskScheduler taskScheduler() {
        ThreadPoolTaskScheduler scheduler = new ThreadPoolTaskScheduler();
        scheduler.setPoolSize(5);
        scheduler.setThreadNamePrefix("aps-scheduler-");
        scheduler.initialize();
        return scheduler;
    }
}
\end{lstlisting}

% ------------------------------------------------------------------------
\subsection{ApiState.java --- Глобальное состояние системы и 7 стадий}
% ------------------------------------------------------------------------

\textbf{Путь:} \texttt{com.smartfactory.state.ApiState}\\
\textbf{Аннотации:} \texttt{@Component}, \texttt{@Data} (Lombok)

Класс реализует потокобезопасный синглтон, хранящий полный слепок состояния цифрового двойника. Помимо исходных общесистемных флагов, класс расширен полями состояния 7 технологических стадий горячего цинкования:

\begin{enumerate}
    \item \textbf{Стадия 1 (\texttt{Stage1Degreasing}):} Щелочное обезжиривание ($T \approx 50^\circ\text{C}$, флаг активности, управление нагревателем ТЭН);
    \item \textbf{Стадия 2 (\texttt{Stage2Rinsing}):} Промывка №1 (каскадный циркуляционный насос);
    \item \textbf{Стадия 3 (\texttt{Stage3Pickling}):} Кислотное травление в соляной кислоте ($T \approx 25^\circ\text{C}$, перемешивающее устройство/барботаж);
    \item \textbf{Стадия 4 (\texttt{Stage4Rinsing}):} Промывка №2 (промывочный насос);
    \item \textbf{Стадия 5 (\texttt{Stage5Fluxing}):} Флюсование ($T \approx 60^\circ\text{C}$, нагреватель);
    \item \textbf{Стадия 6 (\texttt{Stage6DryingChamber}):} Сушильная камера (температура, влажность от физического датчика через UART, вентилятор обдува, индикатор подключения датчика);
    \item \textbf{Стадия 7 (\texttt{Stage7ZincBath}):} Ванна горячего цинкования ($T \approx 450^\circ\text{C}$, расплав цинка, электронагрев).
\end{enumerate}

Метод \texttt{getFactoryStateMap()} формирует структуру данных, идентичную JSON-схеме аппаратного моста на Python.

\begin{lstlisting}[language=Java, caption=Структура ApiState.java и описание 7 стадий]
package com.smartfactory.state;

import lombok.Data;
import org.springframework.stereotype.Component;
import java.util.*;

@Component
@Data
public class ApiState {
    private int temperature = 0;
    private int humidity = 0;
    private boolean doorClosed = true;
    private boolean manualMode = false;
    private String status = "STARTING";
    private String step = "Инициализация";
    private boolean simRunning = false;
    private boolean clientConnected = false;

    // Метрики алгоритма APS
    private String apsPhase = "IDLE";
    private double utilizationCoeff = 0.0;
    private int processingTimeEstimate = 0;
    private boolean canIncludeResult = false;
    private String currentBatch = "";
    private String rejectionCause = "";
    private int recalculationCount = 0;
    private int totalCyclesCompleted = 0;

    // Состояние 7 технологических узлов линии цинкования
    private Stage1Degreasing stage1 = new Stage1Degreasing();
    private Stage2Rinsing stage2 = new Stage2Rinsing();
    private Stage3Pickling stage3 = new Stage3Pickling();
    private Stage4Rinsing stage4 = new Stage4Rinsing();
    private Stage5Fluxing stage5 = new Stage5Fluxing();
    private Stage6DryingChamber stage6 = new Stage6DryingChamber();
    private Stage7ZincBath stage7 = new Stage7ZincBath();

    private List<String> apsHistory = new ArrayList<>();
    private List<LogEntry> systemLog = new ArrayList<>();

    public Map<String, Object> getFactoryStateMap() {
        return Map.of(
            "stage_1_alkaline_degreasing", Map.of("temperature", stage1.getTemperature(), "is_active", stage1.isActive(), "heater_on", stage1.isHeaterOn()),
            "stage_2_rinsing_1", Map.of("is_active", stage2.isActive(), "pump_running", stage2.isPumpRunning()),
            "stage_3_pickling", Map.of("temperature", stage3.getTemperature(), "is_active", stage3.isActive(), "agitator_on", stage3.isAgitatorOn()),
            "stage_4_rinsing_2", Map.of("is_active", stage4.isActive(), "pump_running", stage4.isPumpRunning()),
            "stage_5_fluxing", Map.of("temperature", stage5.getTemperature(), "is_active", stage5.isActive(), "heater_on", stage5.isHeaterOn()),
            "stage_6_drying_chamber", Map.of("temperature_celsius", stage6.getTemperatureCelsius(), "humidity_percent", stage6.getHumidityPercent(), "is_active", stage6.isActive(), "fan_on", stage6.isFanOn(), "sensor_connected", stage6.isSensorConnected()),
            "stage_7_zinc_bath", Map.of("temperature", stage7.getTemperature(), "is_active", stage7.isActive(), "heater_on", stage7.isHeaterOn())
        );
    }

    public synchronized void addApsHistoryEntry(String entry) {
        apsHistory.add(entry);
        if (apsHistory.size() > 20) apsHistory = new ArrayList<>(apsHistory.subList(apsHistory.size() - 20, apsHistory.size()));
    }

    public synchronized void addSystemLog(String source, String message) {
        String time = java.time.LocalTime.now().format(java.time.format.DateTimeFormatter.ofPattern("HH:mm:ss"));
        systemLog.add(new LogEntry(time, source, message));
        if (systemLog.size() > 30) systemLog = new ArrayList<>(systemLog.subList(systemLog.size() - 30, systemLog.size()));
    }

    @Data public static class Stage1Degreasing { private int temperature = 50; private boolean isActive = false; private boolean heaterOn = false; }
    @Data public static class Stage2Rinsing { private boolean isActive = false; private boolean pumpRunning = false; }
    @Data public static class Stage3Pickling { private int temperature = 25; private boolean isActive = false; private boolean agitatorOn = false; }
    @Data public static class Stage4Rinsing { private boolean isActive = false; private boolean pumpRunning = false; }
    @Data public static class Stage5Fluxing { private int temperature = 60; private boolean isActive = false; private boolean heaterOn = false; }
    @Data public static class Stage6DryingChamber { private int temperatureCelsius = 55; private int humidityPercent = 20; private boolean isActive = false; private boolean fanOn = false; private boolean sensorConnected = true; }
    @Data public static class Stage7ZincBath { private int temperature = 450; private boolean isActive = true; private boolean heaterOn = false; }
    @Data @lombok.AllArgsConstructor @lombok.NoArgsConstructor
    public static class LogEntry { private String time; private String source; private String message; }
}
\end{lstlisting}

% ------------------------------------------------------------------------
\subsection{BatchInfo.java --- Модель производственной партии}
% ------------------------------------------------------------------------

\textbf{Путь:} \texttt{com.smartfactory.model.BatchInfo}

Класс описывает партию деталей, включая её идентификатор, наименование, физические параметры (массу, материал) и вычисленный алгоритмом APS динамический приоритет $P$. Содержит фабричный метод \texttt{fromLegacyName(String)} для автоматического синтаксического разбора текстовых заказов оператора с выделением веса и технологического типа сырья.

\begin{lstlisting}[language=Java, caption=Фрагмент модели BatchInfo.java]
package com.smartfactory.model;

import lombok.Data;
import lombok.NoArgsConstructor;
import java.util.ArrayList;
import java.util.List;

@Data
@NoArgsConstructor
public class BatchInfo {
    private String id;
    private String name;
    private String material = "Сталь";
    private double weight = 50.0;
    private int priority = 0;
    private String assignedLineId = "LINE_1";
    private int estimatedTimeSec = 0;
    private BatchStatus status = BatchStatus.WAITING;
    private List<ProcessingStep> steps = new ArrayList<>();

    public enum BatchStatus { WAITING, PRIORITIZED, SCHEDULED, ACTIVE, COMPLETED, REJECTED }

    @Data
    public static class ProcessingStep {
        private int stepNumber;
        private String equipmentId;
        private String operationName;
        private int durationSec;
    }

    public static BatchInfo fromLegacyName(String name) {
        BatchInfo b = new BatchInfo();
        b.setId("BATCH_" + System.currentTimeMillis() % 100000);
        b.setName(name);
        if (name.contains("Медь")) b.setMaterial("Медь");
        else if (name.contains("Алюм")) b.setMaterial("Алюминий");
        else b.setMaterial("Сталь");
        return b;
    }
}
\end{lstlisting}

% ------------------------------------------------------------------------
\subsection{ComplexParameters.java --- Параметры оборудования и расчёт загрузки}
% ------------------------------------------------------------------------

\textbf{Путь:} \texttt{com.smartfactory.model.ComplexParameters}

Инкапсулирует конфигурацию промышленного комплекса, состав производственных линий и перечень единиц оборудования. Реализует математическую модель оценки средней загруженности оборудования $U_{avg}$:
\begin{equation}
U_{avg} = \frac{1}{N} \sum_{i=1}^{N} \frac{L_i}{C_i}
\end{equation}
где $L_i$ --- текущая фактическая загрузка $i$-го агрегата (кг), $C_i$ --- паспортная ёмкость (кг), $N$ --- число задействованных агрегатов с ненулевой ёмкостью. Порог критической перегрузки установлен на уровне $U_{crit} = 0{,}85$ (85\%).

\begin{lstlisting}[language=Java, caption=Расчёт коэффициента загруженности в ComplexParameters.java]
public double getAverageUtilization() {
    double total = 0.0;
    int count = 0;
    for (Equipment eq : equipmentMap.values()) {
        if (eq.getCapacityKg() > 0 && eq.isAvailable()) {
            total += (double) eq.getCurrentLoadKg() / eq.getCapacityKg();
            count++;
        }
    }
    return count > 0 ? total / count : 0.0;
}
\end{lstlisting}

% ------------------------------------------------------------------------
\subsection{SuperpositionData.java --- Файл суперпозиции расписания}
% ------------------------------------------------------------------------

\textbf{Путь:} \texttt{com.smartfactory.model.SuperpositionData}

Реализует концепт суперпозиции: наложение расчётного времени технологических операций партии на временную сетку загрузки оборудования с учётом свободных технологических окон.

\begin{lstlisting}[language=Java, caption=Классы тайм-слота и результата сопоставления в SuperpositionData.java]
@Data
public static class TimeSlot {
    private String equipmentId;
    private int startTimeSec;
    private int endTimeSec;
    private boolean occupied;
    private String batchId;
}

@Data
public static class FitResult {
    private boolean canFit;
    private String rejectionReason;
    private int proposedStartSec;
    private int proposedEndSec;
    private double requiredCapacity;
    private double availableCapacity;
}
\end{lstlisting}

% ------------------------------------------------------------------------
\subsection{ApsAlgorithmService.java --- Реализация 12 блоков алгоритма APS}
% ------------------------------------------------------------------------

\textbf{Путь:} \texttt{com.smartfactory.service.ApsAlgorithmService}\\
\textbf{Размер:} 718 строк\\
\textbf{Зависимости:} \texttt{ApiState}, \texttt{ScheduleService}

Является центральным вычислительным ядром системы. Каждый публичный метод строго соответствует функциональному блоку теоретической блок-схемы алгоритма APS.

\subsubsection{Математическая модель приоритетов}
Приоритет партии $P$ вычисляется как композиция срочности ожидания $U$, коэффициента сырья $M$ и массового фактора $W$:
\begin{equation}
P = U + M + W + P_{manual}
\end{equation}
где $U = \min(50,\, t_{wait} / 10)$, $M \in \{30, 20, 10\}$ (для меди, алюминия и стали соответственно), $W = \min(20,\, m / 5)$. При ручном вклинивании оператора задаётся абсолютный приоритет $P_{manual} = 999$.

\subsubsection{Расчёт времени обработки партии}
Расчётное время $T_{proc}$ учитывает технологическую длительность каждой операции на назначенном оборудовании, скорректированную на температурно-материальный коэффициент $K_m$ и относительную массу заготовки $w$:
\begin{equation}
T_{proc} = \sum_{k=1}^{S} T_{base,k} \cdot K_m \cdot \left(0{,}8 + 0{,}4 \cdot \frac{w}{w_{max}}\right)
\end{equation}

\begin{lstlisting}[language=Java, caption=Фрагменты вычисления приоритета и времени в ApsAlgorithmService.java]
public List<BatchInfo> determineBatchPriority(List<BatchInfo> queue) {
    long now = System.currentTimeMillis();
    for (BatchInfo batch : queue) {
        int urgencyScore = (int) Math.min(50, (now - batch.getCreatedAt()) / 10000);
        int materialScore = switch (batch.getMaterial()) {
            case "Медь" -> 30;
            case "Алюминий" -> 20;
            default -> 10;
        };
        int weightScore = (int) Math.min(20, batch.getWeight() / 5);
        batch.setPriority(urgencyScore + materialScore + weightScore);
    }
    queue.sort((a, b) -> Integer.compare(b.getPriority(), a.getPriority()));
    return queue;
}

public SuperpositionData.FitResult compareWithRealSelection(BatchInfo batch, SuperpositionData supData) {
    SuperpositionData.FitResult result = new SuperpositionData.FitResult();
    double util = complexParams.getAverageUtilization();
    if (util >= 0.85) {
        result.setCanFit(false);
        result.setRejectionReason("Превышен порог загруженности оборудования (" + String.format("%.1f%%", util * 100) + " >= 85%)");
        return result;
    }
    if (batch.getEstimatedTimeSec() > 3600) {
        result.setCanFit(false);
        result.setRejectionReason("Расчётное время обработки превышает допустимый интервал смены");
        return result;
    }
    result.setCanFit(true);
    return result;
}
\end{lstlisting}

% ------------------------------------------------------------------------
\subsection{SimulationService.java --- Симуляция стадий и запуск APS-цикла}
% ------------------------------------------------------------------------

\textbf{Путь:} \texttt{com.smartfactory.service.SimulationService}

Класс объединяет три циклических планируемых процесса:
\begin{enumerate}
    \item \texttt{sensorUpdateLoop()} (каждую 1 секунду) --- сбор и синхронизация показаний датчиков сушильной камеры (DHT) и контроль концевика защитного ограждения;
    \item \texttt{galvanicStagesSimulationLoop()} (каждые 2 секунды) --- математическая симуляция термогидравлических параметров 7 технологических стадий (гистерезис нагревателей, работа циркуляционных насосов, перемешивание барботером, подогрев флюса, расплав цинка при 450$^\circ$C);
    \item \texttt{globalSchedulerLoop()} (каждые 8 секунд) --- последовательное пошаговое выполнение полного цикла блок-схемы APS из 12 этапов с записью результатов в журнал истории.
\end{enumerate}

\begin{lstlisting}[language=Java, caption=Фрагмент SimulationService.java с симуляцией стадий цинкования]
@Scheduled(fixedRate = 2000)
public void galvanicStagesSimulationLoop() {
    if (!apiState.isSimRunning()) return;

    // Узел 1: Щелочное обезжиривание (50-60 C)
    int t1 = 45 + random.nextInt(11);
    apiState.getStage1().setTemperature(t1);
    apiState.getStage1().setHeaterOn(t1 < 50);
    apiState.getStage1().setActive(random.nextBoolean());

    // Узел 2: Промывка 1 (насос)
    boolean pump2 = random.nextBoolean();
    apiState.getStage2().setPumpRunning(pump2);
    apiState.getStage2().setActive(pump2);

    // Узел 3: Травление (20-30 C, мешалка)
    int t3 = 20 + random.nextInt(11);
    apiState.getStage3().setTemperature(t3);
    apiState.getStage3().setAgitatorOn(random.nextBoolean());
    apiState.getStage3().setActive(random.nextBoolean());

    // Узел 4: Промывка 2
    boolean pump4 = random.nextBoolean();
    apiState.getStage4().setPumpRunning(pump4);
    apiState.getStage4().setActive(pump4);

    // Узел 5: Флюс (55-65 C)
    int t5 = 55 + random.nextInt(11);
    apiState.getStage5().setTemperature(t5);
    apiState.getStage5().setHeaterOn(t5 < 60);
    apiState.getStage5().setActive(random.nextBoolean());

    // Узел 7: Ванна цинкования (440-460 C)
    int t7 = 440 + random.nextInt(21);
    apiState.getStage7().setTemperature(t7);
    apiState.getStage7().setHeaterOn(t7 < 448);
    apiState.getStage7().setActive(true);
}
\end{lstlisting}

% ------------------------------------------------------------------------
\subsection{ScheduleService.java --- Потокобезопасная работа с расписанием}
% ------------------------------------------------------------------------

\textbf{Путь:} \texttt{com.smartfactory.service.ScheduleService}

Обеспечивает транзакционное чтение и модификацию файла \texttt{schedule.json}. Защита от состояния гонки между планировщиком APS и внешними REST-запросами клиентов реализована через монитор синхронизации \texttt{fileLock}.

\begin{lstlisting}[language=Java, caption=Фрагмент ScheduleService.java]
@Service
public class ScheduleService {
    private final ObjectMapper mapper = new ObjectMapper();
    private final File scheduleFile = new File("schedule.json");
    private final Object fileLock = new Object();

    public ScheduleData readSchedule() {
        synchronized (fileLock) {
            if (!scheduleFile.exists()) return new ScheduleData();
            try {
                return mapper.readValue(scheduleFile, ScheduleData.class);
            } catch (IOException e) {
                return new ScheduleData();
            }
        }
    }

    public void writeSchedule(ScheduleData data) {
        synchronized (fileLock) {
            try {
                mapper.writerWithDefaultPrettyPrinter().writeValue(scheduleFile, data);
            } catch (IOException e) {
                e.printStackTrace();
            }
        }
    }
}
\end{lstlisting}

% ------------------------------------------------------------------------
\subsection{FactoryController.java --- REST API контроллер}
% ------------------------------------------------------------------------

\textbf{Путь:} \texttt{com.smartfactory.controller.FactoryController}\\
\textbf{Базовый URL:} \texttt{/api}

Предоставляет кросс-доменный HTTP-интерфейс (CORS) для веб-дашборда, мобильного терминала, оркестратора Launcher и Python-моста. Включает специализированные эндпоинты \texttt{/api/factory-status} (двунаправленный обмен телеметрией стадий) и \texttt{/api/health} (liveness-проба готовности ядра).

\begin{lstlisting}[language=Java, caption=Эндпоинты FactoryController.java]
@RestController
@CrossOrigin(origins = "*")
public class FactoryController {
    private final ApiState apiState;
    private final ScheduleService scheduleService;

    @GetMapping("/api/factory-status")
    public Map<String, Object> getFactoryStatus() {
        return Map.of("status", "ok", "data", apiState.getFactoryStateMap());
    }

    @PostMapping("/api/factory-status")
    public Map<String, Object> updateFactoryStatus(@RequestBody Map<String, Object> payload) {
        // Приём телеметрии от Python аппаратного моста
        return Map.of("status", "ok");
    }

    @GetMapping({"/api/health", "/health"})
    public Map<String, Object> healthCheck() {
        return Map.of(
            "status", "UP",
            "opcua_server", "RUNNING",
            "sim_running", apiState.isSimRunning(),
            "manual_mode", apiState.isManualMode()
        );
    }
}
\end{lstlisting}

% ------------------------------------------------------------------------
\subsection{OpcUaServerManager.java --- Менеджер OPC UA сервера}
% ------------------------------------------------------------------------

\textbf{Путь:} \texttt{com.smartfactory.opcua.OpcUaServerManager}

Выполняет жизненный цикл сервера Eclipse Milo на порту 4840. При запуске регистрирует конечную точку \texttt{opc.tcp://0.0.0.0:4840/}, подключает пространство имён \texttt{FactoryNamespace} и запускает периодический цикл синхронизации (каждые 200~мс).

\begin{lstlisting}[language=Java, caption=Синхронизация переменных в OpcUaServerManager.java]
@Scheduled(fixedRate = 200)
public void syncWithOpcUa() {
    if (factoryNamespace == null) return;
    boolean active = factoryNamespace.isClientActive();
    apiState.setClientConnected(active);
    if (!active) return;

    // Считывание команд PyGame-симулятора
    factoryNamespace.readClientValuesToApiState(apiState);
    // Трансляция статусов APS и параметров стадий в узлы OPC UA
    factoryNamespace.writeServerValuesToOpc(apiState);
}
\end{lstlisting}

% ------------------------------------------------------------------------
\subsection{FactoryNamespace.java --- Адресное пространство OPC UA}
% ------------------------------------------------------------------------

\textbf{Путь:} \texttt{com.smartfactory.opcua.FactoryNamespace}\\
\textbf{URI пространства:} \texttt{http://factory-aps.local}

Создаёт и организует иерархическую структуру узлов стандарта IEC 62541. Содержит 8 тематических папок-узлов, полностью покрывающих физическую модель линии цинкования:

\begin{longtable}{|l|l|l|p{6cm}|}
\caption{Структура адресного пространства OPC UA} \\
\hline
\textbf{Узел (Папка)} & \textbf{Переменная} & \textbf{Тип} & \textbf{Назначение} \\
\hline
\endfirsthead
\hline
\textbf{Узел (Папка)} & \textbf{Переменная} & \textbf{Тип} & \textbf{Назначение} \\
\hline
\endhead
\hline
\endfoot
\texttt{1\_Alkaline\_Degreasing} & \texttt{Temperature} & Int32 & Температура щелочного раствора (°C) \\
& \texttt{IsActive} & Boolean & Флаг нахождения детали в ванне \\
& \texttt{HeaterOn} & Boolean & Состояние нагревателя ТЭН \\
\hline
\texttt{2\_Rinsing\_1} & \texttt{IsActive} & Boolean & Флаг промывки детали \\
& \texttt{PumpRunning} & Boolean & Работа циркуляционного насоса \\
\hline
\texttt{3\_Pickling} & \texttt{Temperature} & Int32 & Температура ванны травления (°C) \\
& \texttt{IsActive} & Boolean & Флаг травления детали в кислоте \\
& \texttt{AgitatorOn} & Boolean & Состояние перемешивателя \\
\hline
\texttt{4\_Rinsing\_2} & \texttt{IsActive} & Boolean & Флаг финишной промывки \\
& \texttt{PumpRunning} & Boolean & Работа насоса промывки №2 \\
\hline
\texttt{5\_Fluxing} & \texttt{Temperature} & Int32 & Температура раствора солей флюса (°C) \\
& \texttt{IsActive} & Boolean & Флаг пассивации во флюсе \\
& \texttt{HeaterOn} & Boolean & Подогрев ванны флюсования \\
\hline
\texttt{6\_Drying\_Chamber} & \texttt{TemperatureCelsius} & Int32 & Температура сушки по датчику UART \\
& \texttt{HumidityPercent} & Int32 & Относительная влажность по датчику \\
& \texttt{IsActive} & Boolean & Процесс сушки заготовки \\
& \texttt{FanOn} & Boolean & Включение вентилятора обдува \\
& \texttt{SensorConnected} & Boolean & Статус связи с ESP32 по UART \\
\hline
\texttt{7\_Zinc\_Bath} & \texttt{Temperature} & Int32 & Температура расплава цинка (~450°C) \\
& \texttt{IsActive} & Boolean & Погружение заготовки в расплав \\
& \texttt{HeaterOn} & Boolean & Мощный электроподогрев ванны \\
\hline
\texttt{8\_Global\_Scheduler} & \texttt{SimulationRunning} & Boolean & Пуск/останов симуляции \\
& \texttt{ManualOverride} & Boolean & Флаг ручного вмешательства (кнопка G) \\
& \texttt{SystemStatus} & String & Текущий глобальный статус APS \\
& \texttt{CurrentAlgorithmStep} & String & Название активного шага блок-схемы \\
\hline
\end{longtable}

\begin{lstlisting}[language=Java, caption=Регистрация узлов и запись телеметрии в FactoryNamespace.java]
public void registerNodes() {
    UByte rwAccess = UByte.valueOf(3); // Чтение и запись

    UaFolderNode folder1 = createFolder("1_Alkaline_Degreasing");
    stage1TempNode   = createVariable(folder1.getNodeId(), "1_Alkaline_Degreasing/Temperature", "Temperature", Identifiers.Int32, new Variant(50), rwAccess);
    stage1ActiveNode = createVariable(folder1.getNodeId(), "1_Alkaline_Degreasing/IsActive", "IsActive", Identifiers.Boolean, new Variant(false), rwAccess);
    stage1HeaterNode = createVariable(folder1.getNodeId(), "1_Alkaline_Degreasing/HeaterOn", "HeaterOn", Identifiers.Boolean, new Variant(false), rwAccess);

    // Аналогично регистрируются узлы 2..7

    UaFolderNode sched = createFolder("8_Global_Scheduler");
    simRunningNode     = createVariable(sched.getNodeId(), "SimulationRunning", "SimulationRunning", Identifiers.Boolean, new Variant(false), rwAccess);
    manualOverrideNode = createVariable(sched.getNodeId(), "ManualOverride", "ManualOverride", Identifiers.Boolean, new Variant(false), rwAccess);
    systemStatusNode   = createVariable(sched.getNodeId(), "SystemStatus", "SystemStatus", Identifiers.String, new Variant("WAITING"), rwAccess);
    currentStepNode    = createVariable(sched.getNodeId(), "CurrentAlgorithmStep", "CurrentAlgorithmStep", Identifiers.String, new Variant("Ожидание"), rwAccess);
}

public void writeServerValuesToOpc(ApiState state) {
    if (systemStatusNode != null) systemStatusNode.setValue(new DataValue(new Variant(state.getStatus())));
    if (currentStepNode != null) currentStepNode.setValue(new DataValue(new Variant(state.getStep())));
    if (stage1TempNode != null) stage1TempNode.setValue(new DataValue(new Variant(state.getStage1().getTemperature())));
    if (stage7TempNode != null) stage7TempNode.setValue(new DataValue(new Variant(state.getStage7().getTemperature())));
    // Обновление остальных 6 узлов
}
\end{lstlisting}

% ========================================================================
\newpage
\section{Аппаратный мост и сбор телеметрии (Python / FastAPI / UART)}
% ========================================================================

\subsection{hardware\_bridge.py --- Аппаратно-программный шлюз HIL}

\textbf{Путь:} \texttt{hardware\_bridge.py}\\
\textbf{Язык:} Python 3.10+\\
\textbf{Зависимости:} FastAPI, Uvicorn, PySerial, AsyncIO

Модуль обеспечивает физический уровень сбора данных (Edge/Fog). Выполняет непрерывный неблокирующий опрос последовательного COM-порта (COM20, 115200 бод), к которому подключён микроконтроллер ESP32 с цифровым датчиком DHT11/DHT22.

\subsubsection{Протокол обмена по UART}
ESP32 передаёт пакеты в текстовом формате ASCII, разделённые точкой с запятой:
\begin{verbatim}
temp:24;hum:48\n
\end{verbatim}
Функция \texttt{read\_serial\_data()} парсит значения температуры и относительной влажности, обновляет параметры узла 6 (\texttt{stage\_6\_drying\_chamber}) и переключает признак подключения физического датчика (\texttt{sensor\_connected = True}). При обрыве соединения или отключении кабеля модуль фиксирует ошибку, переводит флаг в \texttt{False} и переключается на встроенный программный генератор.

\subsubsection{Ретрансляция в Java-сервер}
Фоновая корутина \texttt{forward\_to\_java\_server()} с периодичностью 1 раз в секунду выполняет HTTP POST-запрос на адрес \texttt{http://localhost:8080/api/factory-status}, передавая текущее состояние всех узлов в JSON-формате.

\begin{lstlisting}[language=Python, caption=Фрагмент hardware\_bridge.py]
import asyncio, random, urllib.request, json
from contextlib import asynccontextmanager
from fastapi import FastAPI
from fastapi.middleware.cors import CORSMiddleware
try:
    import serial
except ImportError:
    serial = None

COM_PORT = 'COM20'
BAUD_RATE = 115200

factory_state = {
    "stage_1_alkaline_degreasing": {"temperature": 50, "is_active": False, "heater_on": False},
    "stage_2_rinsing_1": {"is_active": False, "pump_running": False},
    "stage_3_pickling": {"temperature": 25, "is_active": False, "agitator_on": False},
    "stage_4_rinsing_2": {"is_active": False, "pump_running": False},
    "stage_5_fluxing": {"temperature": 60, "is_active": False, "heater_on": False},
    "stage_6_drying_chamber": {"temperature_celsius": 0, "humidity_percent": 0, "is_active": False, "fan_on": False, "sensor_connected": False},
    "stage_7_zinc_bath": {"temperature": 450, "is_active": True, "heater_on": False}
}

ser = None

async def read_serial_data():
    global factory_state, ser
    while True:
        try:
            if ser and ser.is_open and ser.in_waiting > 0:
                line = ser.readline().decode('utf-8', errors='ignore').strip().lower()
                if 'temp:' in line and 'hum:' in line:
                    parts = line.split(';')
                    t = int(parts[0].replace('temp:', '').strip())
                    h = int(parts[1].replace('hum:', '').strip())
                    node6 = factory_state["stage_6_drying_chamber"]
                    node6["temperature_celsius"] = t
                    node6["humidity_percent"] = h
                    node6["is_active"] = h > 30
                    node6["fan_on"] = h > 40
                    node6["sensor_connected"] = True
        except Exception as e:
            factory_state["stage_6_drying_chamber"]["sensor_connected"] = False
        await asyncio.sleep(0.1)

async def forward_to_java_server():
    while True:
        try:
            url = "http://localhost:8080/api/factory-status"
            payload = json.dumps({"status": "ok", "data": factory_state}).encode('utf-8')
            req = urllib.request.Request(url, data=payload, headers={'Content-Type': 'application/json'}, method='POST')
            urllib.request.urlopen(req, timeout=1)
        except Exception:
            pass
        await asyncio.sleep(1)
\end{lstlisting}

% ========================================================================
\newpage
\section{Визуальный 2D-симулятор (Python / PyGame)}
% ========================================================================

\subsection{main.py --- Графический SCADA-клиент OPC UA}

\textbf{Путь:} \texttt{pygame-simulator/main.py}\\
\textbf{Зависимости:} PyGame, AsyncUA, AsyncIO, Aiohttp

Модуль реализует цифровой макет физического конвейера линии горячего цинкования. Включает графическую подсистему рендеринга на PyGame и асинхронный клиент OPC UA на базе библиотеки \texttt{asyncua}.

\subsubsection{Двухъярусная топология линии}
Комплекс смоделирован в виде двух связанных линий:
\begin{itemize}
    \item \textbf{Линия 1 (Предварительная подготовка):} Станции 1--4 (Щелочное обезжиривание $\to$ Промывка 1 $\to$ Травление $\to$ Промывка 2);
    \item \textbf{Линия 2 (Нанесение покрытия):} Станции 5--7 (Флюсование $\to$ Сушильная камера с датчиком $\to$ Ванна горячего цинкования).
\end{itemize}

\subsubsection{Автомат состояний мостового манипулятора}
Каждая линия оснащена независимым мостовым краном, логика которого управляется конечным автоматом со следующими состояниями:
\texttt{IDLE} $\to$ \texttt{MOVE\_TO\_PICKUP} $\to$ \texttt{LIFT\_DOWN\_PICKUP} $\to$ \texttt{LIFT\_UP\_PICKUP} $\to$ \texttt{MOVE\_TO\_DROP} $\to$ \texttt{LIFT\_DOWN\_DROP} $\to$ \texttt{LIFT\_UP\_DROP}.

Приоритеты диспетчеризации манипулятора:
1. Выгрузка готовой детали с финальной позиции на выходной конвейер (освобождение линии);
2. Транспортировка деталей между смежными ваннами от конца к началу;
3. Загрузка новой заготовки со входного накопителя на начальную позицию.

\begin{lstlisting}[language=Python, caption=Клиент AsyncUA для 8 узлов в pygame-simulator/main.py]
async def opcua_client_task():
    url = "opc.tcp://localhost:4840/"
    while True:
        client = Client(url=url)
        try:
            await client.connect()
            opc_data["connected"] = True
            idx = await client.get_namespace_index('http://factory-aps.local')

            sched = await client.nodes.objects.get_child([f"{idx}:8_Global_Scheduler"])
            manual_node = await sched.get_child([f"{idx}:ManualOverride"])
            status_node = await sched.get_child([f"{idx}:SystemStatus"])
            step_node   = await sched.get_child([f"{idx}:CurrentAlgorithmStep"])
            sim_node    = await sched.get_child([f"{idx}:SimulationRunning"])

            # Подключение к 7 технологическим папкам-узлам
            f1 = await client.nodes.objects.get_child([f"{idx}:1_Alkaline_Degreasing"])
            s1_t = await f1.get_child([f"{idx}:Temperature"])
            s1_a = await f1.get_child([f"{idx}:IsActive"])
            s1_h = await f1.get_child([f"{idx}:HeaterOn"])

            f6 = await client.nodes.objects.get_child([f"{idx}:6_Drying_Chamber"])
            s6_t = await f6.get_child([f"{idx}:TemperatureCelsius"])
            s6_w = await f6.get_child([f"{idx}:HumidityPercent"])
            s6_s = await f6.get_child([f"{idx}:SensorConnected"])

            while True:
                opc_data["manual_mode"] = await manual_node.get_value()
                opc_data["status"] = await status_node.get_value()
                opc_data["step"] = await step_node.get_value()
                opc_data["sim_running"] = await sim_node.get_value()

                opc_data["stage1"] = {"temp": await s1_t.get_value(), "is_active": await s1_a.get_value(), "heater": await s1_h.get_value()}
                opc_data["stage6"] = {"temp": await s6_t.get_value(), "hum": await s6_w.get_value(), "sensor_connected": await s6_s.get_value()}
                await asyncio.sleep(0.1)
        except Exception:
            opc_data["connected"] = False
            await asyncio.sleep(2)
\end{lstlisting}

% ========================================================================
\newpage
\section{Единый диспетчер запуска (C\# / .NET / Release Launcher)}
% ========================================================================

\subsection{Launcher.cs --- Оркестратор процессов и мониторинг здоровья}

\textbf{Путь:} \texttt{Release/Launcher.cs}\\
\textbf{Язык:} C\# (.NET Framework / Core)\\
\textbf{Размер:} 252 строки

Утилита обеспечивает скоординированный запуск всех компонентов системы в автономной промышленной сборке (standalone release).

\subsubsection{Функции оркестратора}
\begin{enumerate}
    \item \textbf{Валидация окружения:} Проверка установленной виртуальной машины Java (\texttt{java -version});
    \item \textbf{Фоновый запуск ядра:} Старт \texttt{server.jar} с привязкой рабочего каталога;
    \item \textbf{Поллинг готовности:} Циклическая проверка TCP-порта 8080 через \texttt{TcpClient} до подтверждения запуска REST API;
    \item \textbf{Запуск UI:} Автоматическое открытие интерфейса дашборда в браузере по умолчанию и запуск графического симулятора \texttt{simulator.exe};
    \item \textbf{Перехват сигналов завершения:} Регистрация обработчика Win32 API \texttt{SetConsoleCtrlHandler} для гарантированного снятия процессов сервера и симулятора при закрытии консоли.
\end{enumerate}

\begin{lstlisting}[language={[Sharp]C}, caption=Оркестрация и перехват сигналов завершения в Launcher.cs]
using System;
using System.Diagnostics;
using System.IO;
using System.Net.Sockets;
using System.Runtime.InteropServices;
using System.Threading;

class Program {
    private static Process serverProcess = null;
    private static Process simulatorProcess = null;

    [DllImport("Kernel32")]
    private static extern bool SetConsoleCtrlHandler(EventHandler handler, bool add);
    private delegate bool EventHandler(CtrlType sig);
    private enum CtrlType { CTRL_C_EVENT = 0, CTRL_CLOSE_EVENT = 2 }

    static void Cleanup() {
        Console.WriteLine("\n[Launcher] Остановка всех сервисов...");
        try { if (simulatorProcess != null && !simulatorProcess.HasExited) simulatorProcess.Kill(); } catch { }
        try { if (serverProcess != null && !serverProcess.HasExited) serverProcess.Kill(); } catch { }
    }

    static bool CheckPort(int port) {
        try {
            using (TcpClient client = new TcpClient("127.0.0.1", port)) return true;
        } catch { return false; }
    }

    static void Main(string[] args) {
        SetConsoleCtrlHandler((sig) => { Cleanup(); return true; }, true);
        AppDomain.CurrentDomain.ProcessExit += (s, e) => Cleanup();

        string currentDir = AppDomain.CurrentDomain.BaseDirectory;
        string jarPath = Path.Combine(currentDir, "server.jar");
        string simPath = Path.Combine(currentDir, "simulator.exe");

        // Запуск Java Backend
        serverProcess = Process.Start(new ProcessStartInfo {
            FileName = "java", Arguments = "-jar \"" + jarPath + "\"",
            WorkingDirectory = currentDir, UseShellExecute = false
        });

        // Ожидание готовности порта 8080
        int attempts = 0;
        while (attempts++ < 25) {
            if (CheckPort(8080)) break;
            Thread.Sleep(1000);
        }

        Process.Start("http://localhost:8080");
        if (File.Exists(simPath)) {
            simulatorProcess = Process.Start(new ProcessStartInfo { FileName = simPath, WorkingDirectory = currentDir });
        }
    }
}
\end{lstlisting}

% ========================================================================
\newpage
\section{Веб-дашборд (React / TypeScript / Vite)}
% ========================================================================

\subsection{App.tsx --- SCADA-панель цифрового двойника}

\textbf{Путь:} \texttt{src/App.tsx}\\
\textbf{Стек:} React 18, TypeScript, Vite, Tailwind CSS

Интерактивная панель мониторинга для инженера-технолога и оператора АСУ ТП. Осуществляет периодический опрос серверного ядра (каждые 500~мс по эндпоинтам \texttt{/api/aps-state} и \texttt{/api/queue}).

\subsubsection{Функциональные блоки интерфейса}
\begin{enumerate}
    \item \textbf{Телеметрические карточки 7 стадий:} Визуализация статуса активности, температуры с цветовой градацией, состояния насосов, перемешивателей, ТЭНов и датчика DHT11;
    \item \textbf{Интерактивная блок-схема APS:} Динамическая подсветка текущего шага алгоритма синим неоновым свечением, пройденных шагов --- изумрудным цветом;
    \item \textbf{Управление производственной очередью:} Отображение ожидающих партий с приоритетами, ручное добавление/удаление заказов;
    \item \textbf{Пульт ручного вмешательства:} Активация режима \texttt{MANUAL\_INTERVENTION}, кнопки принудительного вклинивания партии в начало расписания или постановки в конец.
\end{enumerate}

\begin{lstlisting}[language=Java, caption=Фрагмент опроса и отображения стадий в App.tsx]
useEffect(() => {
  const fetchState = async () => {
    try {
      const [apsRes, queueRes] = await Promise.all([
        fetch('http://localhost:8080/api/aps-state'),
        fetch('http://localhost:8080/api/queue')
      ]);
      const apsData = await apsRes.json();
      const queueData = await queueRes.json();
      setApsState(apsData);
      setQueueState(queueData);
    } catch (e) {
      console.error('Ошибка поллинга данных:', e);
    }
  };
  const interval = setInterval(fetchState, 500);
  return () => clearInterval(interval);
}, []);
\end{lstlisting}

% ========================================================================
\newpage
\section{Мобильный терминал (Android / Kotlin)}
% ========================================================================

\subsection{Архитектура мобильного клиента}

\textbf{Стек:} Kotlin, Jetpack Compose, Retrofit 2, Moshi, StateFlow

Мобильное приложение представляет собой терминал оператора цеха, реализованный по архитектурному шаблону MVVM (Model-View-ViewModel).

\begin{itemize}
    \item \textbf{\texttt{ApiService.kt}:} Описание REST-интерфейса Retrofit и DTO-моделей данных (\texttt{SystemState}, \texttt{QueueState});
    \item \textbf{\texttt{MainViewModel.kt}:} Содержит реактивные потоки \texttt{StateFlow} и корутины \texttt{viewModelScope}, выполняющие периодический опрос сервера;
    \item \textbf{\texttt{MainActivity.kt}:} Декларативный пользовательский интерфейс на Jetpack Compose с Material Design 3.
\end{itemize}

\begin{lstlisting}[language=Java, caption=Сетевой интерфейс Retrofit в ApiService.kt]
package com.example.api

import retrofit2.http.GET
import retrofit2.http.POST
import retrofit2.http.Body

interface ApiService {
    @GET("/api/state")
    async fun getState(): SystemState

    @GET("/api/queue")
    async fun getQueue(): QueueState

    @POST("/api/queue")
    async fun addTask(@Body request: TaskRequest): TaskResponse

    @POST("/api/toggle-manual")
    async fun toggleManual(): SystemState
}
\end{lstlisting}

% ========================================================================
\newpage
\section{Вспомогательные файлы и конфигурации}
% ========================================================================

\subsection{schedule.json --- Персистентное расписание}

Формат файла расписания на диске:
\begin{lstlisting}[language=Java, caption=Структура schedule.json]
{
  "queue": [
    {
      "id": "BATCH_1721000000_001",
      "name": "Партия Сталь А-12",
      "material": "Сталь",
      "weight": 45.0,
      "priority": 78,
      "status": "WAITING",
      "steps": []
    }
  ],
  "current_tasks": []
}
\end{lstlisting}

% ========================================================================
\newpage
\section{Заключение}
% ========================================================================

Разработанный программно-аппаратный комплекс включает \textbf{18 программных модулей}, распределённых по 5 ключевым подсистемам: серверному ядру, аппаратному шлюзу HIL, графическому симулятору SCADA, оркестратору запуска и клиентским пользовательским интерфейсам. Общий объём разработанного исходного кода превышает \textbf{3500 строк}.

Ключевым достоинством архитектуры является её модульность, строгая типизация адресного пространства OPC UA и прямое взаимно-однозначное соответствие между теоретической блок-схемой алгоритма APS и сервисным слоем \texttt{ApsAlgorithmService.java}. Система обеспечивает бесперебойную синхронизацию физических данных датчиков с цифровым двойником и высокую гибкость при переходе в режим ручной диспетчеризации.

\end{document}
